\section{Partial differential equations}

\subsection{First order hyperbolic PDEs: The advection equation}

A few explicit first order numerical methods for solving the advection equation 
$\frac{\partial u}{\partial t} = -v \frac{\partial u}{\partial x}$
are given below: 
\begin{eqnarray}
\mathrm{Forward~time~centered~space~(FTCS):} & u^{n+1}_j = \frac{-v \Delta t}{2\Delta x} \, \Big[{u^n_{j+1} - u^n_{j-1}}\Big] + u^n_j \\
\mathrm{Upwind~method~(for~} v > 0): & u^{n+1}_j = \frac{-v \Delta t}{\Delta x} \, \Big[{u^n_{j} - u^n_{j-1}}\Big] + u^n_j \\
\mathrm{Downwind~method~(for~} v < 0): & u^{n+1}_j = \frac{-v \Delta t}{\Delta x} \, \Big[{u^n_{j+1} - u^n_{j}}\Big] + u^n_j~.
\end{eqnarray}
Above, $u^{n+1}_j = u(t_n, x_j)$, where $x_j = x_0 + j \Delta x$ and $t_n = t_0 + n \Delta t$.
\subsubsection{Problems:}
\begin{enumerate}
\item Solve the advection equation with $v = 1$ and initial condition $u(t=0,x) := \exp[-(x-x_0)^2]$ (Gaussian pulse) where $x \in [0, 25]$ and $x_0 = 5$ using the FTCS scheme, using the following condition for the left boundary $u(x=0,t) = 0$. Is the evolution stable?  
\item Repeat the calculation with upwind methods using Courant factor $\lambda := |v\Delta t|/\Delta x = 1/2, 1, 2$. Which of the evolutions are stable? 
\item Repeat the calculation using the periodic boundary condition $u(x=0, t) = u(x=L, t)$ where $x=L$ corresponds to the right boundary. 
\item Plot the order of convergence of $u(x,t)$ as a function of $x$ at $t = 10$ (use $\lambda = 1)$.   
\end{enumerate}

\subsection{Second order hyperbolic PDEs: The wave equation}

A simple explicit method for solving the wave equation $\frac{\partial^2 u}{\partial t^2} - v^2 \frac{\partial^2 u}{\partial x^2} = 0$ is given below: 
\begin{eqnarray}
\mathrm{For~the~first~time~step:} & u^{n+1}_j = \frac{\lambda^2}{2} \left( u^n_{j+1} + u^n_{j-1} \right) + \left(1 - \lambda^2 \right) u^n_j + \Delta t \frac{\partial u}{\partial t} \big |^n_j\\ 
\mathrm{All~other~time~steps:} & u^{n+1}_j = \lambda^2 \left( u^n_{j+1} + u^n_{j-1} \right) + 2 \left(1 - \lambda^2 \right) u^n_j - u^{n-1}_j\, , 
\end{eqnarray}
where $\lambda = v\Delta t/\Delta x$ is the Courant factor.

\subsubsection{Problems:}
\begin{enumerate}
\item Solve the wave equation with $v = 1$ and initial conditions $u(t=0,x) := \sin \pi x$ and $\frac{\partial u}{\partial t}(t=0,x) := 0$, and Dirichlet type boundary conditions $u(t,x=0) = u(t,x=L) = 0 $ over the domain $x \in (0,L)$ and $t \in (0,T)$. Choose $L = 4$ and $T = 20$. Be sure to choose $\lambda \leq 1$. This shows oscillations on a string whose end points are fixed.  
\end{enumerate}

\subsection{Second order parabolic PDEs: The diffusion equation}

The FTSC method provides a simple explicit scheme for solving the diffusion equation $\frac{\partial u}{\partial t} - k \frac{\partial^2 u}{\partial x^2} = 0$:
\begin{eqnarray}
u_j^{n+1} = \gamma \left(u^n_{j+1} + u^n_{j-1} \right) + (1 - 2 \gamma) u^n_j
\end{eqnarray}
where $\gamma = \frac{k\Delta t}{\Delta x^2}$. The FTSC scheme can produce stable evolutions for $\gamma \leq 0.5$. 

\subsubsection{Problems}
\begin{enumerate}
\item Solve the diffusion equation with $k = 0.5$ over the domain $x \in [0, 2],~t \in [0, 10]$ with the following initial and boundary conditions: 
\begin{eqnarray}
\label{eq:heat_ini_cond}
u (t=0, x) & = & 
\begin{cases}
100 x,~\mathrm{for}~ 0 \leq x \leq 1 \\
200-100 x,~\mathrm{for}~ 1 < x \leq 2 \\
\end{cases}\\
u(t, x=0) & = & 50 \\
u(t, x=2) & = & 10 
\end{eqnarray}
This shows the evolution of the temperature profile (initial profile given by Eq.~\ref{eq:heat_ini_cond}) of a one-dimensional object connected to hot baths at both ends, maintained at temperatures 50 K and 10 K. 
\end{enumerate}

\subsection{Second order elliptic PDEs: The Laplace equation}

The Laplace equation $\frac{\partial^2 u}{\partial x^2} + \frac{\partial^2 u}{\partial y^2} = 0$ in two dimensions can be solved using the successive relaxation method. This explicit method involves successively iterating the following equation 
\begin{equation}
u_{i,j} = \frac{u_{i-1,j} + u_{i+1,j} + u_{i,j-1} + u_{i,j+1}}{4}.
\end{equation}
Above, $u_{i,j} = u(x_i, y_j)$, where $x_i = x_0 + i h$ and $y_j = y_0 + j h$. The iteration has to be repeated until the difference of $u_{i,j}$ between two successive iterations is less than the required tolerance. That is, $\epsilon := \sum_{i,j} \, |u_{i,j}^{n} - u_{i,j}^{n-1}| \leq \mathrm{tol}$.

If a von Neumann type boundary condition $\frac{\partial u}{\partial x} = A(y)$ is specified at $x = 0$, then the $i=0$ grid points can be evaluated using the following equation  
\begin{equation}
u_{0,j} = \frac{2 u_{1,j} - 2 h A_j + u_{0,j-1} + u_{0,j+1}}{4}.
\end{equation}

\subsubsection{Problems}

\begin{enumerate}
\item Solve the Laplace equation over a $100 \times 100$ grid  with the following Dirichlet type boundary conditions:
\begin{equation}
u(x = 0, y) = u(x, y=0) = u(x, y=L) = 0, ~~ u(x = L, y) = 100.  
\label{eq:laplace_bc}
\end{equation}
Plot $u(x,y)$ once the system has reached equilibrium state $\epsilon \lesssim 10^{-6}$). This shows the steady-state temperature distribution on a two-dimensional plate which is kept at a fixed temperature of zero degree at all boundaries, except at the right boundary (which is kept at 100 K). 
\item Plot the error $\epsilon$ as a function of iteration. One should see the error monotonically decreasing. 
\item Repeat the calculation, except that the boundary condition at $x = 0$ in Eq.\,(\ref{eq:laplace_bc}) is replaced by $\frac{\partial u}{\partial x} = 10$.
\end{enumerate}

\section{Fourier and spectral methods}
\subsection{Fast Fourier transform (FFT)}
\subsubsection{Problems:}
\begin{enumerate}
\item Hercules X-1 is a high-mass X-ray binary (HMXB) system having a magnetized spinning neutron star which happens to be an X-ray pulsar. Here~\cite{rxte-data} you are given the data of an RXTE~\cite{rxte} observation after cleaning and pre-processing. The file contains two columns, (i) time (in seconds) and (ii) count-rate of X-ray photons (i.e., number of photons detected per unit time). Plot the count-rate as function of time. This is called lightcurve in X-ray astronomy. Compute the FFT of the given time-series and plot its absolute value against the frequency. Are you able to find any periodic signal(s)? The lowest frequency signal corresponds to the spin frequency of the neutron star.  What is the spin period of this pulsar? Can you identify other periodic signals and how they are related to the lowest frequency signal?~\footnote{This problem and data are graciously provided by Dr. Arunava Mukherjee (IUCAA).}
\end{enumerate}

\subsection{Power spectrum estimation using FFT}
\subsubsection{Problems:}
\begin{enumerate}
\item A sample data set from the LIGO gravitational-wave observatory can be downloaded from here~\cite{ligo-data}. This contains 256 seconds of LIGO data from 2005, sampled at a rate of 4096 Hz. Compute the power spectral density of the data using Welch's modified periodogram method. 
\end{enumerate}

\subsection{Time-frequency signal detection methods}
\subsubsection{Problems}
\begin{enumerate}
\item 4U 1636--536 is a low-mass X-ray binary (LMXB) system consisting of a rapidly spinning weakly magnetized neutron star which does not exhibit any coherent X-ray pulsation like Hercules X-1. However, X-ray radiation from this source shows signals of high frequency (500--1000 Hz) quasi-periodic oscillations (QPOs). Here~\cite{rxte-data2} you are given another data of RXTE observation. This contains the arrival times of photons. Identify the frequency span of this QPO using a time-frequency spectrogram or a PSD. {Note: Firstly, you need to sample the data with a fixed sampling rate, say 4096 Hz}.  
\end{enumerate}

\subsection{Computing correlations using FFT: Matched filtering}
In the case a known signal $h(t)$ buried in stationary Gaussian, white noise, the optimal technique for 
signal extraction is the \emph{matched filtering}, which involves cross-correlating the data 
with a \emph{template} of the signal. The correlation function between two time series $x(t)$ and $\hat h(t)$ (both assumed to be real-valued) for a time shift 
$\tau$ is defined as:
\begin{equation}
\label{eq:corrFn}
R(\tau) = \int_{-\infty}^{\infty}x(t)\,\hat{h}(t+\tau)\,dt. 
\end{equation}
Above, $\hat{h(t)} := h(t) / ||h||$, where the norm $||h||$
of the template is defined by $$||h||^2 = \int_{0}^{t_c} \, |h(t)|^2/\sigma^2 \, dt,$$ where 
$\sigma^2$ is the variance of the noise. 
The optimal signal-to-noise ratio (SNR) is obtained when the template exactly matches with the signal. i.e., $\mathrm{SNR}_\mathrm{opt} = ||h||$. 
If the SNR is greater than a predetermined threshold (which corresponds to an acceptably small false alarm 
probability), a detection can be claimed. 

\subsubsection{Problems}
\begin{enumerate}
\item You are given a time-series data set $d(t)$ here~\cite{mock-gw-data}. This contains a simulated gravitational-wave signal from a black hole binary buried in zero-mean, Gaussian white noise $n(t)$ with standard deviation $\sigma = 10^{-21}$. i.e.,
\begin{equation}
d(t) = n(t) + m \, h_+(t) / D_L,
\end{equation}
where $h_+(t)$ is given by Eq.~(\ref{eq:GWpolzn}), $m = m_1 + m_2$ is the total mass of the binary, and $D_L$ is the (unknown) luminosity distance to the binary.  Detect the location of the signal in the data by maximizing the correlation of the waveform templates computed in Eq.~(\ref{eq:GWpolzn}) with the data.: 
\begin{equation}
R_\mathrm{max} = \mathrm{max}~_{m,\mu,\tau}~ R(\tau),
\end{equation}
where $R(\tau)$ is given by Eq.~(\ref{eq:corrFn}). We have the prior information that the parameters of the signal are in the following range: $5 < m/M_\odot < 15$, $0.2 < \mu/m < 0.25$. Estimate the parameters $m,\mu$ of the signal (parameters that maximize the correlation). 

\end{enumerate}

